\section{Human Psychometrics}

\subsection{Goal}

The goal of the human workstream was to establish whether the sparse ARC human dataset is good enough to serve as a benchmark at all, and if so, what kind of benchmark it is. A second goal was to calibrate how much agreement we should expect between humans and models given the dataset's own internal noise.

\subsection{Data}

The human attempt log contains 4,681 attempts across 509 sessions, 442 observed task IDs, and 502 task-pair rows. Coverage is sparse and opportunistic rather than exam-like: only 40.4\% of ARC-AGI-2 public task pairs appear in the observed table, and the implied session-by-item matrix density is about 1.83\%. Public Train is easier than Public Eval, but Public Eval is the main overlap region for cross-system comparison.

\subsection{Methods}

The human pipeline combines descriptive solve-rate summaries with person-plus-item logistic modeling, split-half simulation, and later latent task estimation in the cross-ARC package. The key methodological move is to treat split-half human-versus-human correlations as the reference distribution for how large any human-versus-model correlation can realistically be on this sparse benchmark. This keeps later model comparisons from being judged against an unrealistic ceiling.

\subsection{Experiment}

The human line effectively ran three experiments. First, it tested whether the sparse human matrix still supports meaningful psychometric estimation. Second, it measured human split-half reproducibility on the public-eval overlap used for cross-system comparison. Third, it compared human item difficulty to several model-side profiles: the average model, the best single model, and a per-pair oracle profile.

\subsection{Results}

The basic result is that the human benchmark is noisy but usable. The held-in diagnostics for the person-plus-item model are strong enough to support latent estimation, and the cross-ARC package later shows that latent task scores are more stable than raw task means. On the public-eval overlap, the human split-half median Pearson correlation is about 0.399. Human versus average-model correlation is about 0.402, which lands near the center of the human split-half distribution. Human versus best single model is lower, around 0.276, and human versus per-pair oracle sits in between, around 0.343.

This leads to a balanced reading. The strongest version of ``models are random with respect to humans'' is not supported, because the average-model profile clearly aligns with human item difficulty. But the strongest version of ``leaderboard success means human-like ordering'' is also not supported, because even a strong single model does not line up with humans as well as one split-half of humans lines up with the other.

The human side also has its own distinctive burden signature. Human difficulty is much more associated with duration than with raw board size, and within-task pair-level heterogeneity is substantial. These features become crucial later when the paper asks why solver complexity predicts LLM difficulty more strongly than human difficulty.

\begin{table}[t]
\centering
\small
\caption{Human-versus-model alignment in the public-eval overlap, benchmarked against human split-half reliability.}
\begin{tabular}{p{2.6in}p{1.0in}p{1.2in}}
\toprule
\textbf{Series} & \textbf{Pearson $r$} & \textbf{Percentile vs split-half} \\
\midrule
Human split-half median & 0.399 & 0.500 \\
Human vs average model & 0.402 & 0.523 \\
Human vs best single model & 0.276 & 0.023 \\
Human vs per-pair oracle & 0.343 & 0.179 \\
\bottomrule
\end{tabular}
\end{table}

\begin{figure}[H]
    \centering
    \includegraphics[width=0.92\textwidth]{human_split_half_stability.png}
    \caption{Raw solve rates versus latent task estimates under split-half resampling in the cross-ARC package. The main point is that latent task estimates are more stable than naive raw means in every benchmark slice.}
\end{figure}

\subsection{Why this section matters for the rest of the paper}

The human workstream is what forces the paper to distinguish shared difficulty from shared burden. Without it, ARC could be narrated as a clean machine psychometric instrument. With it, the paper has to deal with sparse exposure, moderate reliability, duration, and pair-level heterogeneity. Those are not side issues; they are the empirical reasons the full manuscript ends up arguing that humans and models overlap on one difficulty axis while diverging on what makes a task burdensome.
